% ch01 —— §1 Introduction and Summary（原文 wcf-paper.tex 199–555 行）
\section{引言与摘要}
\label{sec:Introduction}

本文的主题是 $d=4$、$\CN=2$ 超对称量子场论中 BPS 态简并度的一个\textbf{穿墙公式}（wall-crossing formula，下称 WCF）。我们的约定与相关定义的汇总见第 2 节。电荷为 $\gamma$ 的 BPS 态空间 $\CH_{\gamma,\mathrm{BPS}}$，指的是单粒子 Hilbert 空间中的这样一个空间：其中各态的电磁荷为 $\gamma$，且使 BPS 界 $M \geq \abs{Z_\gamma(u)}$ 取等号。这里 $u$ 记矢量多重态模空间 $\vm$（即真空模空间的库仑分支）中的一点。

对 $d=4$、$\CN=2$ 超对称而言，唯一可用的指标是第二螺旋度超迹：
\begin{equation} \label{eq:index}
\degen(\gamma;u) := -\half\,{\Tr}_{\CH_{\mathrm{BPS},\gamma}}(-1)^{2J_3}(2J_3)^2,
\end{equation}
其中 $J_3$ 是有质小群转动子群的任一生成元。人们早已知道，这类指标一般并不独立于 $u$，而只是分段常值 \cite{Cecotti:1992qh}。事实上，$\degen(\gamma;u)$ 会在边缘稳定墙处跳变，在那里 $\gamma=\gamma_1+\gamma_2$ 且 $\arg Z_{\gamma_1}(u)=\arg Z_{\gamma_2}(u)$。这一事实在 \SW 理论的发展中起过重要作用 \cite{Seiberg:1994rs,Seiberg:1994aj}。

近些年来，人们对这一指标如何依赖于 $u$ 开始有了更为系统的理解。当衰变 $\gamma\to\gamma_1+\gamma_2$ 中至少有一个组分是原始的（primitive）时，\cite{Denef:2007vg} 给出了跨过边缘稳定墙时变化量 $\Delta\degen$ 的公式。这些原始与半原始穿墙公式，是由基于超引力多中心解的物理图像导出的 \cite{Denef:2000nb,Denef:2000ar}。然而，当两个组分都带非原始电荷时，\cite{Denef:2007vg} 的方法便难以施展。

Kontsevich 与 Soibelman \cite{ks1} 提出了一个适用于一切可能衰变情形的、关于 $\Delta\degen$ 的漂亮穿墙公式。我们将在第 2.2 节回顾他们的公式，下文有时称之为 KS 公式。

一方面，Kontsevich--Soibelman 的「Donaldson--Thomas 不变量」$\hat\degen(\gamma;u)$ 与物理上所关心的 $\degen(\gamma;u)$ 显然并不相同，他们导出公式所用的技巧也似乎与标准的物理考量相去较远。另一方面，与 \cite{Denef:2007vg} 中半原始穿墙公式的表述相比，他们的 WCF 牵涉到一些引人注目的新概念。特别是，WCF 是用环面的一族辛同胚的某个乘积来表达的（见下文 (2.18)），该乘积依赖于诸 $\hat\degen(\gamma;u)$，从而\emph{先验地}依赖于 $u$。WCF 的陈述是：这个乘积实际上与 $u$ 无关。这反过来就决定了 $\hat\degen(\gamma;u)$ 对 $u$ 的依赖方式。这一进展提出了 KS 公式的物理推导与物理诠释问题，也预示着其中含有某些本质性的新物理思想。事实果然如此。

在本文中，我们针对 $d=4$、$\CN=2$ 的\emph{场论}给出 KS 公式的物理解释与证明。推广到超引力是一个有趣而重要的问题，留待将来。

以下是主要思想与基本策略的梗概。考虑该规范理论处在空间 $\IR^3\times S^1$ 上，其中 $S^1$ 的半径为 $R$。低能下该理论由以超凯勒流形 $(\CM,g)$ 为靶空间的 $d=3$ sigma 模型描述。该 sigma 模型受到 BPS 瞬子的修正：所谓 BPS 瞬子，就是 $d=4$ 理论中某个 BPS 粒子的世界线绕 $S^1$ 缠绕一圈所形成的组态。于是，把度量 $g$ 在大 $R$ 处展开，便可读出 BPS 粒子的简并度 $\degen(\gamma;u)$。这立刻引出一个困惑：我们知道诸 $\degen(\gamma;u)$ 是不连续的，而 $g$ 却应当是连续的！度量的连续性基于这样一条物理原则（它在 \cite{Seiberg:1994rs,Seiberg:1994aj} 中至关重要）：低能有效场论拉氏量中的奇异性，只能来自在特殊模处出现的无质量粒子——这些粒子本就不该被从有效理论中积分掉。

物理上，这一困惑的化解类似于 \cite{GarciaEtxebarria:2007zv} 中最近讨论的情形。精确的度量 $g$ 确实是光滑的，但它既受单粒子态修正，也受多粒子态修正。当某个粒子衰变时，其单瞬子贡献的消失，正好被它的衰变产物所带来的多瞬子贡献中的不连续性所补偿。同样地，消失的 $n$-瞬子贡献被 $m>n$ 的 $m$-瞬子贡献中的不连续性所补偿。更精确地说，$n$-瞬子修正具有如下形式
\[
\sum \prod_{i=1}^{n}\degen(\gamma_i;u)\,F^{(n)}_{(\gamma_i)}(u),
\]
其中求和遍历电荷的一切 $n$ 元组 $\{\gamma_i\}$，而 $F^{(n)}$ 本质上是 $R$ 与诸 $Z_{\gamma_i}$ 的普适函数。跨过墙时，每个 $\degen(\gamma;u)$ 都有一个不连续量，它正比于满足 $\sum\gamma_j=\gamma$ 的诸 $\degen(\gamma_j;u)$ 之乘积的和；与此同时，诸函数 $F^{(n)}$ 也带有不连续量，正比于下标 $n'<n$ 的诸 $F^{(n')}$。我们将会看到，Kontsevich--Soibelman 穿墙公式所表达的，正是这一连串抵消的自洽性。

因此，理解 WCF 的主要技术障碍就在于：对 BPS 瞬子给 $g$ 带来的修正给出一种高效的描述。超凯勒度量是个复杂的对象，若去研究它各个分量的修正，几乎难有进展；一般说来，也没有像凯勒势那样简单的可加对象可用。为克服这一困难，我们借用扭量理论的一些想法。回忆：一个超凯勒流形相对于整整一个 $\IC\IP^1$ 的复结构都是复辛的。基本想法是：研究 $g$ 等价于研究 $\CM$ 上的全纯达布坐标——只要我们同时考虑\emph{全部}这些复结构。

在本文主体中，我们假定 Kontsevich--Soibelman 穿墙公式对 $\degen(\gamma;u)$ 成立。在这一假定下，我们构造 $\CM$ 上的度量，办法是给出一族标准函数 $\CX_\gamma(u,\theta;\zeta)$，它们定义在 $\CM\times\IC^\times$ 上，并由电磁荷 $\gamma$ 标记。这里 $(u,\theta)$ 标定 $\CM$ 中一点，参数 $\zeta$ 标记 $\CM$ 上的复结构。每个 $\CX_\gamma$ 关于 $\zeta$ 是\emph{分段}全纯的；BPS 瞬子的效应就是在 $\zeta$ 平面上沿某些射线 $\ell$ 制造出 $\CX_\gamma(u,\theta;\zeta)$ 的不连续性。这些不连续性被等同于 Kontsevich 与 Soibelman 引入的那些辛同胚。在这一进路下，度量的连续性成为 WCF 的推论。最后一节我们把论证倒过来运行：利用超对称规范理论的一般原理，推出度量 $g$ 的若干性质，它们足以证明 WCF。

\subsection*{小结}

我们首先在第 2 节回顾 $d=4$、$\CN=2$ 规范理论的 \SW 解以及 Kontsevich--Soibelman 穿墙公式，然后讨论该理论在 $\IR^3\times S^1$ 上的表述。它是一个映入流形 $\CM$ 的 sigma 模型；$\CM$ 在拓扑上是 $d=4$ 模空间 $\vm$ 上的 \SW 环面丛，并装备一个超凯勒度量 $g$。该度量依赖于 $S^1$ 的半径 $R$。当 $R\to\infty$ 时它趋于一个简单形式，该形式可由 $d=4$ 理论的朴素量纲约化得到；我们称这个简单度量为 $g^{\sfm}$（sf 取自 ``semiflat''）。

在第 3 节我们解释构造超凯勒度量的「扭量」方法：给定 $\CM$ 上一族函数 $\CX_\gamma(u,\theta;\zeta)$，它们随 $\zeta\in\IC^\times$ 全纯变化并满足若干附加条件，就存在一个超凯勒度量，使得这些 $\CX_\gamma(u,\theta;\zeta)$ 成为它的全纯达布坐标。特别地，我们给出对应于半平坦度量 $g^{\sfm}$ 的函数 $\CX^{\sfm}_\gamma(u,\theta;\zeta)$。

有了这些背景，我们就可以考虑瞬子修正。第 4 节从最简单的情形开始：$U(1)$ 规范理论耦合一个电荷为 $q$ 的物质超多重态。在这一理论中，修正后的度量 $g$ 是精确已知的 \cite{Ooguri:1996me,Seiberg:1996ns}。我们说明如何通过纳入瞬子修正、把函数 $\CX^{\sfm}_\gamma(u,\theta;\zeta)$ 改造为新的 $\CX_\gamma(u,\theta;\zeta)$ 来得到这一修正度量。在这一构造中，Kontsevich--Soibelman 公式的构件已经现身：我们的 $\CX_\gamma(u,\theta;\zeta)$ 自然而然地带有 $\zeta$ 平面中的不连续性，而这些不连续性恰恰就是对应于电荷 $\pm q$ 的初等 Kontsevich--Soibelman 辛同胚。

随后，第 5 节转向更有意思的情形：存在多种 BPS 瞬子修正，它们来自 $d=4$ 中相互非局域的 BPS 粒子。此时我们为 $\CX_\gamma(u,\theta;\zeta)$ 找到一个自然的拟设（ansatz）：本质上，只要求每个 BPS 粒子各自独立地贡献一个不连续性，与我们为单个粒子找到的那一个同型。这一不连续性最自然地落在 $\zeta$ 平面中由该 BPS 粒子中心荷的相位所决定的一条射线上。相互非局域粒子的 Kontsevich--Soibelman 因子并不对易；但一般说来这不构成问题，因为这些粒子的中心荷互不共线，它们的不连续性于是出现在 $\zeta$ 平面中彼此不同的射线上。射线的分离恰在边缘稳定墙上消失；在那里，诸不连续性堆积成 Kontsevich--Soibelman 因子的乘积。WCF 说的正是：从墙的两侧分别趋近时，这一乘积相同。这一要求对我们至关重要：它蕴含着由 $\CX_\gamma(\zeta)$ 构造出的度量是连续的。

更确切地说，为确定 $\CX_\gamma(u,\theta;\zeta)$，我们既要指定它们在 $\zeta$ 平面中的不连续性，又要指定它们在 $\zeta\to0,\infty$ 时的渐近行为。换言之，我们提出一个无穷维的「黎曼--希尔伯特问题」，其解就是诸 $\CX_\gamma(u,\theta;\zeta)$。我们并不精确地构造它的解；我们采用的策略与 Cecotti 和 Vafa 所用者极为接近——他们在研究 $d=2$、$\CN=(2,2)$ 超对称理论时遇到过类似的（但为有限维的）黎曼--希尔伯特问题 \cite{Cecotti:1993rm}。把他们论证稍作变形，就足以证明我们的问题至少在 $R$ 充分大时存在解。事实上，在大 $R$ 极限下，所需的 $\CX_\gamma$ 可以逐次逼近得到：第零阶近似就是 $\CX^{\sfm}_\gamma$，第 $n$ 阶近似则把至多 $n$ 个瞬子的多瞬子效应纳入进来。

在构造出诸 $\CX_\gamma(u,\theta;\zeta)$、从而构造出度量 $g$ 之后，我们核验 $g$ 具有若干基于一般场论理由所能预期的性质；它通过了全部检验，因此我们论证它应当就是 $\CM$ 上正确的物理度量，这推广了 \cite{Seiberg:1996nz} 中的一个类似论证。

如前所述，我们对 $\CX_\gamma(u,\theta;\zeta)$ 的构造与 $d=2$ 情形 \cite{Cecotti:1993rm} 中出现的构造有着惊人的相似。在那一情形，BPS 畴壁简并度的穿墙公式是利用 $d=2$ 理论真空丛上平坦的「$tt^*$ 联络」证明的。该联络的两个分量给出表达 $d=2$ 理论 $R$ 对称性与标度不变性的微分方程。我们的构造同样可以表述为：在 $\CM_u$（即 $\CM$ 的环面纤维）上实解析函数的无穷维丛中，$\vm\times\IC\IP^1\times\IR_+$ 上的一个平坦联络 $\CA$。黎曼--希尔伯特构造保证了这一 $\CA$ 的存在性。每个 $\CX_\gamma$ 定义一个平坦截面。特别地，这一平坦性给出关于 $\CX_\gamma$ 对 $\zeta$ 与 $R$ 的依赖性的—对微分方程，它们在物理上源自 $d=4$ 理论的 $R$ 对称性反常与标度不变性。

正如我们在第 7 节所述，这个类 $tt^*$ 的平坦联络，仅利用超对称规范理论的一般原理就能被发现。而且，它仅仅\emph{存在}这一事实，就足以先验地证成我们对度量的拟设。特别地，穿墙公式原本是在该拟设内部作为一条自洽性要求而出现的，现在它可以理解为由 $\CA$ 构造出的一个「等单值形变」的存在性。这就给出了穿墙公式的一个物理证明。

为方便起见，本文大部分篇幅使用不包含味对称信息的简单形式的穿墙公式，相应地也把所有味质量取为零。第 6 节将说明如何恢复味荷与质量的信息。

我们附上若干包含更多细节的附录。附录 A 说明如何直接验证穿墙公式在纯 $SU(2)$ 理论情形给出正确的 BPS 简并度。附录 B 描述 $\CM$ 上的柯西--黎曼方程，其表述与我们对超凯勒度量的构造、以及 \cite{Cecotti:1991me} 的 $tt^*$ 方程相衔接。附录 C 给出从我们的黎曼--希尔伯特问题中提取度量大 $R$ 修正所需的渐近分析。附录 D 讨论从黎曼--希尔伯特问题的解中提取微分方程的若干细节。最后，附录 E 解释我们的主要结果之一，即方程 (5.13)，与热力学贝特拟设之间的一个有趣联系。这一联系还有更多可说的，但留待另一次。

文中有若干小节专门讨论与 KS 公式中一个细微却重要的符号相关的整体性问题。初读时跳过这些讨论是合理的。选择这样做的读者，尽可允许自己在正文中把 $T$ 与 $\tilde T$、以及 $\CM$ 与 $\widetilde{\CM}$ 混为一谈。

\subsection*{讨论}

我们来点评几个特别有意思的地方。

\begin{itemize}
\item 物理上，我们对 $\CM$ 上度量的构造，相当于给出了一条在 $d=4$ 中把相互非局域的粒子「积分掉」的规则。这个问题\emph{先验地}看来很难，因为找不到任何一个使所有粒子都带电荷的对偶标架，从而很难写出包含全部相关场的拉氏量。这里我们绕开了这一困难。

\item 我们对度量的构造用的是它的扭量描述。扭量空间出现的最自然的物理语境是射影超空间 \cite{Karlhede:1984vr,Lindstrom:1987ks,Lindstrom:2008gs}，其中参数 $\zeta$ 是一个玻色型超空间坐标。$g$ 的修正只来自 BPS 瞬子、并且定域在 $\zeta$ 平面中特定的射线上，这一事实在射影超空间的语言中应当有自然的解释。

\item Kontsevich--Soibelman WCF 的灵感之一，是他们早年的工作 \cite{MR2181810}：他们通过「修正」半平坦 K3 上的函数层，给出了 K3 曲面上全纯函数层的一个构造。修正是用某些辛同胚的乘积来表述的，与穿墙公式中出现的那些类似。这一构造与我们的构造密切相关，只需把 K3 换成 $\CM$。我们工作中的关键新要素是同时考虑所有复结构，从而引入参数 $\zeta\in\IC\IP^1$；一旦这样做了，我们就能写出那个至关重要的黎曼--希尔伯特问题。这个想法或许在原来的 K3 语境中也有用。

\item $g$ 的多瞬子展开是一些基本构件之和，各构件以 BPS 简并度 $\degen(\gamma;u)$ 的乘积为权重。这些基本构件在边缘稳定墙处有精细的不连续行为，它与 $\degen(\gamma;u)$ 的跳变相互配合，使 $g$ 关于 $u$ 连续。这一切令人想起 Joyce 关于穿墙的近期工作 \cite{Joyce:2006pf}。此外，Joyce 的工作被 Bridgeland 与 Toledano Laredo 在 \cite{stab-stokes} 中解释为 $\IC\IP^1$ 上某个联络的等单值形变；该联络与我们此处考虑的联络有些相像，但形式略有不同：它只在 $t=0$ 处有非正则奇点，而我们的联络在 $\zeta=0$ 与 $\zeta=\infty$ 两处都有非正则奇点。我们的联络有一个有趣的标度极限：$R\to0$、$\zeta\to0$ 而 $\zeta/R=t$ 固定，它把我们的联络带到 \cite{stab-stokes} 中联络的形式（尽管结构群不同）。这一极限保留了关于 BPS 简并度及其穿墙的信息。看看是否存在某种意义把我们的联络与 \cite{stab-stokes} 中的联系起来，将会很有意思。

\item 我们的讨论中研究的是定义在矢量多重态模空间 $\vm$ 上的结构。然而 Kontsevich--Soibelman 与 Joyce 都把他们的不变量表述在一个更大的空间上，即「Bridgeland 稳定性条件」的空间 \cite{MR2373143}。我们尚不理解，把这些构造推广到该更大空间时的含义为何。

\item Kontsevich--Soibelman 所表述的穿墙公式，不仅对 $\CN=2$ 场论、而且对超引力都有意义，事实上后者正是 \cite{Denef:2007vg} 的主要关注点。此时理论在 $\IR^3\times S^1$ 上的模空间 $\CM$ 是四元数凯勒流形而非超凯勒流形。尽管如此，经由适当修改，我们的多数考量在该语境中似乎仍然说得通。例如，Hitchin 定理被 LeBrun 定理取代，后者用全纯接触结构来刻画四元数凯勒流形的扭量空间。特别地，仍然有「全纯」函数 $\CX_\gamma(x,\zeta)$ 的自然概念（即 $\CM$ 的扭量空间上的全纯函数），而 $\CX^{\sfm}_\gamma$ 的四元数凯勒类比已在 \cite{Neitzke:2007ke} 中求出。我们预期 $\CM$ 上受瞬子修正的度量 $g$，可以用与本文平行的方法得到：为 $\CX_\gamma$ 写出黎曼--希尔伯特问题，用 $\CX^{\sfm}_\gamma$ 定渐近行为，用 Kontsevich--Soibelman 因子定不连续性。需要克服的一个重要困难是：在引力情形，简并度 $\degen(\gamma;u)$ 随 $\gamma$ 增长得极快，这使得迭代解 $\CX_\gamma$ 的收敛性在此情形下不那么显然。与超凯勒情形一样，WCF 应当作为保证 $g$ 光滑的一条自洽性条件而出现。

\item 丛 $\CM\to\vm$ 的超凯勒几何与 \cite{Cecotti:1991me,Cecotti:1992qh,Cecotti:1993rm} 的 $tt^*$ 几何之间的类似是惊人的：两种结构非常相像，尽管一个关乎 $d=4$ 的场论，另一个关乎 $d=2$。两者之间有没有直接的联系？一种可能是仅通过紧致化把它们联系起来，例如在 $S^2$ 上紧致化。$S^2$ 上 $U(1)$ 磁通的不同取值于是对应于 $d=2$ 理论的不同真空，而 $d=4$ 理论的 BPS 态则可被等同于 $d=2$ 中在这些真空之间插值的畴壁。\footnote{这一图像是 Cumrun Vafa 向我们倡导的。}相关的想法此前已在文献中出现过——特别见 \cite{Losev:1997tp,Losev:1999nt,Ooguri:2005vr}。另见 \cite{Dorey:1999zk,Ritz:2006rt}，其中给出了 $d=2$ 与 $d=4$ 中 BPS 谱之间一种略有不同的联系。

\item 包含半平坦几何在内的一类超凯勒流形的无穷小形变，最近在 \cite{Alexandrov:2008ds} 中得到了研究。用他们的语言描述半平坦几何的首阶修正将会很有意思。我们的方程 (4.33) 与他们的 (3.38) 相像，只要适当选取 $H$ 与积分围道。
\end{itemize}

\begin{yizhu}
本章是全书译稿的样板，以下术语与记号在此一并交代，后文沿用。
\begin{enumerate}[leftmargin=2em,itemsep=2pt,topsep=2pt]
\item \textbf{墙与穿墙}：wall 译「墙」，wall-crossing 与 wall-crossing formula 依数译既有译法作「穿墙」「穿墙公式」（不用「墙穿越」）；wall of marginal stability 作「边缘稳定墙」（不用「边际稳定」）。BPS 简并度指标 $\degen(\gamma;u)$ 即原文的 \verb|\degen|，为第二螺旋度超迹，沿用 $\Omega$ 这一符号；Kontsevich--Soibelman 的不变量则记为 $\hat\degen(\gamma;u)$，正文中称 $\hat\Omega$ 或「KS 不变量」。
\item \textbf{情形与语境}：原文以「$d=4$ 理论放在 $\IR^3\times S^1$ 上」的方式过渡到三维，故 instanton 一词在此指「BPS 粒子沿 $S^1$ 缠绕」所形成的组态，译「BPS 瞬子」，与四维瞬子区别开；multi-instanton 作「多瞬子」，$n$-instanton 作「$n$-瞬子」。
\item \textbf{人名与缩写}：Kontsevich--Soibelman 首次出现用全名，其后作「KS 公式」「KS 因子」「KS 辛同胚」；Seiberg--Witten 作「\SW 理论」；人名之间用 en dash（--），与中文连接号区分。
\item \textbf{交叉引用}：本章提及的公式 (2.18)、(4.33)、(5.13) 与各节号均按原著排版给出；合稿时统一改为自动编号，以免因增删文字而错位。
\end{enumerate}
\end{yizhu}
